\documentclass[11pt]{article}
\usepackage[margin=1in]{geometry}
\usepackage{amsmath,amssymb,amsthm,booktabs,needspace}
\usepackage[hidelinks]{hyperref}
\setlength{\emergencystretch}{2em}
\newtheorem{theorem}{Theorem}
\newtheorem{lemma}[theorem]{Lemma}
\newtheorem{corollary}[theorem]{Corollary}
\newcommand{\mr}{\operatorname{mr}}
\newcommand{\rank}{\operatorname{rank}}
\newcommand{\Z}{\mathbb Z}
\newcommand{\F}{\mathbb F}
\title{Additivity and limits of simple branched-cover rank bounds}
\author{Internal research working proof}
\date{17 November 2026}
\begin{document}
\maketitle

For a classical knot $K$, let $Y_\phi$ be the connected
simple $r$-sheeted branched cover associated to a homomorphism
$\phi:\pi_1(S^3\setminus K)\twoheadrightarrow S_r$ taking
meridians to transpositions. Define
\[
 C(K)=\max_{r,\phi}\{r-1+\rank\pi_1(Y_\phi)\},\qquad
 A_k(K)=\max_{r,\phi}\{r-1+b_1(Y_\phi;k)\},
\]
where $r\ge2$ and $k$ is a fixed field. The maxima are
well-defined: the double cover always exists, and the
retained bound
\begin{equation}\label{eq:bound}
                r-1+\rank\pi_1(Y_\phi)\le\mr(K)\le b(K)
\end{equation}
forces $r\le b(K)+1$. There are finitely many homomorphisms
to each finite $S_r$. In particular $1\le A_k(K)\le C(K)\le\mr(K)$.

\begin{theorem}\label{thm:add}
For all classical knots $K_1,\ldots,K_t$,
\[
 C(\mathop{\#}_{i=1}^t K_i)=\sum_i C(K_i)-(t-1),\qquad
 A_k(\mathop{\#}_{i=1}^t K_i)=\sum_i A_k(K_i)-(t-1).
\]
The same field $k$ is used in every term of the second identity.
\end{theorem}

The proof classifies the restrictions of every simple
coloring of a connected sum, then counts the spheres in
the lifted decomposition. The ordinary-rank step imports
the classical Grushko theorem explicitly. No general
meridional-rank additivity theorem is assumed or proved.
For example, the known knot $T(4,7)$ has
\[
                     C(T(4,7))=3<4=\mr(T(4,7)).
\]
Thus even the strongest rank bound in this class of covers
does not solve full MRC. The uniform statements below are
written working proofs, unformalized and externally unreviewed;
worldwide priority is not claimed.

\section{The retained meridional cover bound}

For completeness, if the knot group is generated by $d$
meridians, its lifted graph has $r$ vertices. Each generator
exchanges two sheets, and all other lifted edges are fixed-sheet
meridian loops. Branched filling kills those loops and the
two-edge meridian-square cycle. The exchanged edges can
therefore be identified with opposite orientations, leaving
$d$ edges on $r$ vertices. The graph is connected; killing
a spanning tree leaves $d-r+1$ generators for the cover group.
The filling loops used here are valid for every generating
meridian, since they are conjugates of the same boundary
meridians. This gives \eqref{eq:bound} for arbitrary
meridional generating families. Replacing group rank by
first homology gives the scalar bound in any field.

\section{Normal generation leaves one active set of sheets}

\begin{lemma}\label{lem:orbit}
If a knot group maps to $S_R$ and a meridian maps to a
transposition $\tau=(0\,1)$, its image is the full symmetric
group on one subset $S$ containing $\{0,1\}$, fixing every
sheet outside $S$ pointwise.
\end{lemma}
\begin{proof}
A knot group is normally generated by one meridian: all
Wirtinger generators are conjugate meridians. Thus its
image $G$ is generated by its conjugates of $\tau$.
Let $S$ be the $G$-orbit of 0, which also contains 1.
Every such conjugate has both endpoints in $S$ and fixes
the complement. Hence so does all of $G$.
The graph on $S$ with these transpositions as edges is
connected, because $S$ is an orbit. Transpositions along
a connected graph generate the full symmetric group on
its vertices. Therefore $G=S_S$.
\end{proof}

This lemma uses the knot hypothesis. A link group need
not be normally generated by one chosen component meridian,
so the same conclusion is not asserted for arbitrary links.

Write the group of the connected sum as the amalgam of
the summand knot groups identifying their meridians.
Given a simple $R$-coloring of the sum, conjugate the common
meridian image to $(0\,1)$. Lemma~\ref{lem:orbit} gives
subsets $S_i$, each containing $\{0,1\}$, on which the
restricted image is $S_{r_i}$ with $r_i=|S_i|$.
Their union is all $R$ sheets, since the full coloring is
transitive. The restrictions give connected simple covers
$Y_i$ of the summands on their active sheets. On each of
the other $R-r_i$ sheets, the summand cover is a separate
copy of $S^3$.

Conversely, start with any connected simple covers of the
summands. Relabel their meridian transpositions to $(0\,1)$
and embed their active sheet sets in a common set. As long
as the sets contain $\{0,1\}$ and cover the common set,
the homomorphisms agree on the amalgamated meridian.
They define a connected simple coloring of the sum.
Indeed, their transposition graphs are connected through
the common pair, so their combined image is $S_R$.

\section{The lifted sphere decomposition}

\begin{theorem}\label{thm:glue}
For active sheet sets $S_i$ as above, the cover $Y$ of the
connected sum satisfies
\[
 Y\cong\mathop{\#}_{i=1}^t Y_i\ \#\
       \mathop{\#}^{\beta}(S^1\times S^2),\qquad
 \beta=\sum_i(r_i-2)-(R-2).
\]
In particular $\beta\ge0$.
\end{theorem}
\begin{proof}
Choose the decomposing spheres for the connected sum so
that every neck meets the knot twice on arcs whose meridian
image is $(0\,1)$. Each sphere lifts to $R-1$ spheres:
one double cover branched over its two intersection points,
and $R-2$ degree-one copies. All these surfaces are $S^2$.

The cover of a ball containing a single trivial proper arc
likewise consists of $R-1$ balls. On the moving pair the
map is the product of $z\mapsto z^2$ on a disk with an
interval; all fixed-sheet covers are single balls.
Removing the small summing balls from each summand cover
therefore punctures its components by balls. Its capped
components are $Y_i$ and $R-r_i$ copies of $S^3$.

Form a graph with these capped components as vertices and
the lifted summing spheres as edges. It is connected, because
the full monodromy is transitive. There are
\[
 V=\sum_i(1+R-r_i),\qquad E=(t-1)(R-1)
\]
vertices and edges. Gluing along a spanning tree forms the
connected sum of all the capped vertex manifolds. Each
remaining edge joins two spherical holes of the same connected
oriented manifold, adding an $S^1\times S^2$ summand.
The number of these edges is
\[
 E-V+1=\sum_i(r_i-2)-(R-2)=\beta.
\]
The $S^3$ vertex summands have no effect. This proves the formula.
No irreducibility or geometric recognition of any $Y_i$ is required.
\end{proof}

For equal sheet sets of size $r$, this gives
$\beta=(t-1)(r-2)$. For sheet sets whose only common sheets
are the moving pair and whose other sheets are disjoint,
$R=2+\sum_i(r_i-2)$ and $\beta=0$.
Different sheet choices trade cover degree for free factors.

\section{Proof of additivity and sharpness consequences}

By van Kampen, Theorem~\ref{thm:glue} gives
\[
 \pi_1(Y)\cong\left(*_i\pi_1(Y_i)\right)*F_\beta.
\]
The classical Grushko theorem implies
\[
 \rank\pi_1(Y)=\sum_i\rank\pi_1(Y_i)+\beta.
\]
Here all groups are finitely generated. The precise
free-group lifting theorem is quoted in Stallings
\cite[Theorem 1.B.5]{stallings}; applying it to a free group
on a minimal generating set of the displayed free product
gives the lower rank inequality, and the union of factor
generating sets gives the upper inequality. This is an
explicit classical import; the original 1965 proof is not
independently re-certified here.

Consequently every connected coloring of the sum satisfies
\begin{align*}
 R-1+\rank\pi_1(Y)
 &=\sum_i\bigl(r_i-1+\rank\pi_1(Y_i)\bigr)-(t-1).
\end{align*}
Since every coloring restricts as in Lemma~\ref{lem:orbit},
this is at most $\sum_i C(K_i)-(t-1)$.
Conversely, choose a maximizing cover of each summand and
embed the active sheet sets with only the moving pair in
common. The converse construction above realizes that sum
of bounds. This proves the first equality in
Theorem~\ref{thm:add}.

For a fixed field $k$, first homology of the displayed free
product has dimension $\sum_i b_1(Y_i;k)+\beta$.
The identical maximization argument proves the second
equality. This scalar assertion does not require Grushko's
ordinary-rank theorem. Maximizing separately over different
fields on the summands would be a different operation;
the fixed-field hypothesis is retained.

\begin{corollary}\label{cor:sharp}
If $C(K_i)=\mr(K_i)$ for every summand, then
\[
 \mr(\mathop{\#}_i K_i)=\sum_i\mr(K_i)-(t-1).
\]
If instead one has matching cover lower and bridge upper
certificates $C(K_i)=b(K_i)=b_i$, then
\[
 \mr(\mathop{\#}_i K_i)=b(\mathop{\#}_i K_i)
                       =\sum_i b_i-(t-1).
\]
\end{corollary}
\begin{proof}
The additive cover bound gives the lower inequality.
Meridional generating families of the factors can be
conjugated so that one chosen generator is the common
meridian; their union then uses one fewer generator at
each join. This supplies the matching meridional upper
bound. Joining bridge positions supplies the stated
bridge upper bound as well.
\end{proof}

Schubert's bridge additivity, in the primary modern proof
\cite[Theorem 1]{schultens}, further shows that the
nonnegative defects $b-C$ and $b-A_k$ are additive under
connected sum. Thus the class $C=b$ is closed under
connected sum and taking connected-sum factors.
This does not assert such a characterization for the
class $C=\mr$ without bridge information.

The earlier four periodic families have $C=4$: their
three-block finite action induces a connected simple
three-cover whose group maps onto the nonabelian block
stabilizer. Fixed-sheet meridian loops and meridian squares
act trivially, so this map descends through branched filling.
The upper-four bound then forces that cover rank to be exactly two.
Therefore any connected sum of $t$ members of those families
and $s$ trefoils, with $t+s\ge1$, has
\[
                              \mr=b=3t+s+1.
\]
The trefoil double cover has group $C_3$ and hence $C=2$.
The assertion is a uniform extension of the retained
certificates; membership in prior established families and
worldwide novelty remain unestablished.

\section{An exact limitation: the known knot T(4,7)}

Use the standard torus presentation and meridian convention
already checked in the retained affine proof, following
Cornwell \cite[Section 5.1]{cornwell}:
\[
 \Gamma=\langle x,y\mid x^4=y^7\rangle,\qquad m=x^{-1}y^2.
\]
The retained binary four-dimensional rank-one representation
and the four-strand braid give $\mr=b=4$.
For any simple connected $r$-coloring, \eqref{eq:bound}
gives $r\le5$. If $3\le r\le5$, the central element
$x^4=y^7$ maps to the trivial center of $S_r$.
Then $y^7=1$, but $S_r$ has no nonidentity element of
order seven. Thus $y$ maps to 1 and the image is cyclic,
a contradiction. Only the double cover remains.

In the orbifold group impose $t=m$, $t^2=1$, so $x=y^2t$.
The double-cover group is its index-two kernel, with $y$
even and $t$ odd. Schreier rewriting with transversal
$\{1,t\}$ and $z=tyt$ gives
\[
 G=\langle y,z\mid y^2z^2y^2z^2=y^7,
                         \ z^2y^2z^2y^2=z^7\rangle.
\]
In particular $\rank G\le2$. Its abelian exponent matrix
$\left(\begin{smallmatrix}-3&4\\4&-3\end{smallmatrix}\right)$
has first Smith factor one and determinant $-7$, so
$H_1(G;\Z)=\Z/7$.

Nonabelianity is certified over $\F_2$ by
\[
 X=\begin{pmatrix}1&0&0\\1&1&0\\0&1&1\end{pmatrix},\qquad
 Y=\begin{pmatrix}0&1&0\\0&0&1\\1&1&0\end{pmatrix},\qquad
 M=X^{-1}Y^2,\quad Z=MYM.
\]
Direct evaluation gives $X^4=Y^7=M^2=I$; $Y,Z$ satisfy
both displayed cover relators and do not commute.
Thus $\rank G=2$, proving
\[
 C(T(4,7))=3,\qquad
 A_k(T(4,7))=\begin{cases}2&\operatorname{char}k=7,\\
                         1&\text{otherwise}.
                  \end{cases}
\]
This is a limitation of these lower bounds on an already
settled MRC knot, not a counterexample to MRC.

For $J_n=\#^nT(4,7)$, Theorem~\ref{thm:add} now gives
\[
 C(J_n)=2n+1,\qquad A_k(J_n)=1+n\,[\operatorname{char}k=7].
\]
Schubert additivity gives $b(J_n)=3n+1$; MRC for these
graph-manifold-exterior knots is already covered by
Boileau--Dutra--Jang--Weidmann \cite[Corollary 1]{bdjw}.
Their published 2017 theorem is distinct from the later
2022 Whitehead preprint with the retained proof gap.
The deficits of the cover-rank and scalar tests therefore
grow without bound even on known MRC cases.

There is also no onto $S_3$ transposition coloring of any
$J_n$: the orbit lemma would restrict it to such a coloring
of a factor, or to degree-two supports all equal to the same
moving pair. The first is impossible, and the second cannot
generate $S_3$. By the retained rank-two equivalence, the
rank-two-form part of the original binary quadratic criterion
has no representation on $J_n$ in any dimension. No exclusion
of all higher-rank forms on these connected sums is asserted.

\section{Finite controls, prior work, and assurance}

The frozen eight-control plan constructs literal connected-sum
diagrams and sheet assignments in degrees two through five.
For two rank-two seed covers with homology $\Z/37$ and
$\Z/35$, sharing all three sheets gives
\[
 Y_1\#Y_2\#(S^1\times S^2),\quad\rank\pi_1=5,
 \quad H_1=\Z\oplus\Z/1295.
\]
Sharing only the moving pair gives a four-cover with
$Y_1\#Y_2$, rank four and finite homology $\Z/1295$.
Both lower bounds are seven. For three seed factors,
the analogous degrees three and five give ranks eight
and six, respectively, and the same lower bound ten.
Mixed trefoil and unknot controls check degree and unit
boundaries.

The independent verifier traces arcs between undercrossings,
checks the planar knot walk, reconstructs the lifted words
in reverse order, and replays every unimodular Smith step.
It separately checks the two-state Schreier words and the
binary torus matrices on all eight vectors. Four corrupted
palette/free-factor/Smith/torus inputs are rejected.
These computations check the specified inputs and homology;
the all-knots additivity claim requires the written orbit,
sphere-gluing and classical rank arguments above.

Conditional meridional additivity for knots detected by
Coxeter quotients is already noted in Baader--Blair--Kjuchukova
\cite[Introduction]{bbk}. No general priority claim for
conditional additivity or branched-cover rank methods is made.
The field-scalar invariant is the retained symmetric root-lift
bound, so its additivity records an exact connection to that
part of the unpublished programme. The separate original sixteen
prime targets are settled in the existing BJW18 graph-cover MRC family;
see \path{ORIGINAL_SIXTEEN_KNOWN_GRAPH_COVER_FAMILY.tex}. The other
panel names and full classical MRC remain open.

\clearpage
\begin{thebibliography}{9}
\bibitem{stallings} J. Stallings, \emph{Group Theory and
Three-Dimensional Manifolds}, Yale University Press (1971),
Theorem 1.B.5, printed p.~6.
\url{https://webhomes.maths.ed.ac.uk/~v1ranick/papers/stallings1.pdf}.
Original topological proof: Math. Z. 90 (1965), 1--8,
\url{https://doi.org/10.1007/BF01112046}; publisher metadata
checked, full original proof not audited.
\bibitem{schultens} J. Schultens, Additivity of bridge numbers
of knots, Math. Proc. Cambridge Philos. Soc. 135 (2003),
539--544. Theorem 1; credit to Schubert (1954).
\url{https://doi.org/10.1017/S0305004103006832}.
Author manuscript:
\url{https://www.math.ucdavis.edu/~jcs/pubs/addbri15.pdf}.
\bibitem{cornwell} C. Cornwell, Knot contact homology and
representations of knot groups, J. Topology 7 (2014),
1221--1242. Section 5.1.
\url{https://arxiv.org/abs/1303.4943}.
\bibitem{bdjw} M. Boileau, E. Dutra, Y. Jang and R. Weidmann,
Meridional rank of knots whose exterior is a graph manifold,
Topology Appl. 228 (2017), 458--485. Corollary 1.
\url{https://arxiv.org/abs/1608.01570}.
\bibitem{bbk} S. Baader, R. Blair and A. Kjuchukova,
Coxeter groups and meridional rank of links,
Math. Ann. 379 (2021), 1533--1551.
\url{https://doi.org/10.1007/s00208-020-02124-z}.
\end{thebibliography}
\end{document}
